%============================================================
% BAB 2 — hasil konversi awal dari DOCX lama
%============================================================
\chapter{KAJIAN PUSTAKA}

Bab ini menguraikan landasan konseptual dan teoretis yang digunakan sebagai dasar untuk membangun penelitian deteksi potensi \emph{fraud} atau \emph{upcoding} pada klaim JKN/BPJS dalam sistem INA-CBG. Pembahasan dimulai dari deskripsi problem domain pada sistem klaim kesehatan nasional, kemudian dilanjutkan dengan teori penunjang yang menjelaskan dasar analisis berbasis data, serta diakhiri dengan penelitian terkait untuk menempatkan posisi dan orisinalitas penelitian ini.

\section{DESKRIPSI PERMASALAHAN}\label{deskripsi-permasalahan}

Permasalahan utama dalam proyek akhir ini adalah bagaimana mendeteksi potensi \emph{fraud} atau \emph{upcoding} pada klaim JKN/BPJS dalam sistem INA-CBG secara lebih terstruktur, terukur, dan dapat dijelaskan. Pada domain klaim kesehatan, deteksi tidak dapat dipahami hanya sebagai klasifikasi teknis antara klaim normal dan klaim bermasalah, tetapi harus dilihat dalam konteks sistem pembayaran, proses pengkodean, verifikasi klaim, dan keterbatasan ketersediaan label \emph{fraud} aktual. Oleh karena itu, deskripsi permasalahan pada bagian ini menempatkan klaim kesehatan sebagai objek utama yang dianalisis, sekaligus menegaskan bahwa keluaran penelitian diarahkan untuk mendukung prioritas verifikasi, bukan untuk membuktikan \emph{fraud} aktual secara hukum atau administratif.

\subsection{JKN/BPJS, FKRTL, dan Mekanisme Klaim INA-CBG}\label{jknbpjs-fkrtl-dan-mekanisme-klaim-ina-cbg}

Program Jaminan Kesehatan Nasional merupakan skema pembiayaan kesehatan yang dikelola oleh BPJS Kesehatan untuk menjamin akses pelayanan kesehatan bagi peserta. Pada fasilitas kesehatan rujukan tingkat lanjutan, pengajuan klaim dilakukan melalui mekanisme \emph{Indonesian Case-Based Groups} atau INA-CBG, yaitu sistem pembayaran prospektif berbasis kelompok kasus yang menetapkan besaran pembayaran berdasarkan kombinasi diagnosis, prosedur, dan karakteristik klinis tertentu. Permenkes Nomor 27 Tahun 2014 menjelaskan bahwa sistem INA-CBG merupakan bagian dari penyelenggaraan \emph{casemix} dalam program JKN, sedangkan petunjuk teknis \emph{e-claim} menunjukkan bahwa proses klaim dibentuk melalui entri data administratif dan klinis, proses \emph{grouping}, finalisasi klaim, dan pengiriman data klaim ke sistem verifikasi \citep{kemenkes2014}, \citep{kemenkes2016}.

Dalam konteks penelitian ini, FKRTL menjadi lokasi yang sangat relevan karena pada level inilah nilai klaim sangat dipengaruhi oleh ketepatan dokumentasi diagnosis, prosedur, kelas perawatan, dan tingkat keparahan kasus. Dengan demikian, klaim di FKRTL bukan hanya catatan biaya, melainkan representasi interaksi antara kondisi pasien, keputusan klinis, pengkodean rekam medis, dan aturan pembayaran. Karakter inilah yang menjadikan data klaim JKN/BPJS pada FKRTL penting untuk dianalisis ketika tujuan penelitian diarahkan pada deteksi potensi \emph{fraud} atau \emph{upcoding}.

\subsection{\texorpdfstring{Diagnosis, Prosedur, \emph{Severity Level}, \emph{Grouping}, dan Tarif Klaim}{Diagnosis, Prosedur, Severity Level, Grouping, dan Tarif Klaim}}\label{diagnosis-prosedur-severity-level-grouping-dan-tarif-klaim}

Pada sistem INA-CBG, nilai klaim ditentukan melalui proses \emph{grouping} yang mempertimbangkan diagnosis utama, diagnosis sekunder, prosedur, jenis pelayanan, dan \emph{severity level}. Struktur ini menunjukkan bahwa perubahan pada data klinis dan administratif dapat memengaruhi kelompok kasus yang dihasilkan beserta tarif yang melekat padanya. Petunjuk teknis e-claim dan dokumen INA-CBG menunjukkan bahwa proses entri klaim tidak berhenti pada pencatatan penyakit, tetapi berujung pada pembentukan kode kasus dan tarif pembayaran yang menjadi dasar pengajuan klaim kepada BPJS Kesehatan \citep{kemenkes2016}.

Hubungan antara komponen klaim dan struktur tarif inilah yang membuat sistem pembayaran berbasis kelompok kasus memiliki konsekuensi insentif. Chalkley dkk. menunjukkan bahwa perubahan harga dalam sistem pembayaran \emph{diagnosis-related groups} di Indonesia berkaitan dengan perubahan perilaku \emph{coding} rumah sakit, meskipun besarannya terbatas \citep{chalkley2022}. Temuan tersebut penting karena memperlihatkan bahwa struktur pembayaran tidak netral terhadap perilaku pencatatan klinis. Crespin dkk. (2024) juga menegaskan bahwa sistem pembayaran berbasis diagnosis dapat mendorong pengodean pasien ke tingkat intensitas atau keparahan yang lebih tinggi \citep{crespin2024}. Oleh sebab itu, diagnosis, prosedur, \emph{severity level}, \emph{grouping}, dan tarif klaim harus dipahami sebagai rangkaian unsur yang saling terkait dalam membentuk risiko ketidakwajaran klaim.

\subsection{\texorpdfstring{\emph{Upcoding} sebagai Bagian dari \emph{Fraud} atau \emph{Abuse Provider}}{Upcoding sebagai Bagian dari Fraud atau Abuse Provider}}\label{upcoding-sebagai-bagian-dari-fraud-atau-abuse-provider}

Dalam literatur \emph{fraud} kesehatan, \emph{upcoding} dipahami sebagai praktik penggunaan kode diagnosis atau prosedur yang menghasilkan pembayaran lebih tinggi daripada yang seharusnya. Villegas-Ortega dkk. menempatkan praktik semacam ini dalam spektrum \emph{health insurance fraud} sebagai bentuk misrepresentasi untuk memperoleh manfaat yang tidak sah \citep{villegas2021}. Maryana dan Aspan menjelaskan bahwa upcoding terjadi ketika diagnosis atau prosedur diubah menjadi kode yang bertarif lebih tinggi daripada kondisi klinis sebenarnya \citep{maryana2024}. Dalam konteks JKN/BPJS, praktik tersebut relevan karena tarif INA-CBG sangat bergantung pada hasil \emph{grouping}, sehingga perubahan kode dapat berdampak langsung pada besaran klaim.

\emph{Upcoding} pada \emph{provider} juga dapat dibaca sebagai bagian dari \emph{moral hazard} atau \emph{abuse} dalam sistem pembiayaan kesehatan. Syafrawati dkk. menunjukkan bahwa \emph{moral hazard provider} dalam skema asuransi kesehatan sosial dapat muncul dalam bentuk \emph{upcoding}, \emph{readmission}, dan \emph{unnecessary admission} \citep{syafrawati2023}. Muliarta dkk. serta Nurlianti dkk. juga menegaskan bahwa \emph{fraud} dalam JKN mencakup berbagai bentuk manipulasi administratif dan klinis yang dapat merugikan sistem pembiayaan kesehatan \citep{nurlianti2025}, \citep{muliarta2023}. Dengan demikian, \emph{upcoding} tidak berdiri sendiri sebagai kesalahan teknis, tetapi perlu ditempatkan dalam konteks perilaku organisasi dan insentif finansial pada sistem klaim kesehatan.

\subsection{\texorpdfstring{\emph{Fraud} Aktual, \emph{Error Coding}, dan Indikasi Risiko}{Fraud Aktual, Error Coding, dan Indikasi Risiko}}\label{fraud-aktual-error-coding-dan-indikasi-risiko}

Salah satu tantangan utama dalam domain ini adalah membedakan antara \emph{fraud} aktual, \emph{error coding}, dan indikasi risiko. \emph{Fraud} aktual mensyaratkan unsur kesengajaan dan biasanya memerlukan verifikasi administratif, klinis, atau audit untuk dibuktikan. Sebaliknya, \emph{error coding} dapat terjadi karena ketidaktepatan dokumentasi, perbedaan interpretasi aturan, keterbatasan \emph{coder}, atau ketidaklengkapan rekam medis. Sugiarti dkk. menunjukkan bahwa ketidaktepatan kode diagnosis dan ketidaksesuaian \emph{clinical pathway} dapat menghasilkan tarif klaim yang lebih tinggi, tetapi \emph{upcoding} tidak selalu dapat langsung disebut \emph{fraud} apabila unsur kesengajaannya belum dapat dibuktikan \citep{sugiarti2021}.

Perbedaan tersebut memiliki implikasi metodologis yang penting. Model analitik yang dibangun dari data klaim hanya membaca pola administrasi, biaya, histori layanan, dan karakteristik klinis yang terekam di dalam data. Karena itu, keluaran model lebih tepat diposisikan sebagai indikasi risiko atau sinyal ketidakwajaran klaim, bukan sebagai penetapan \emph{fraud} aktual. Posisi ini konsisten dengan kebutuhan verifikasi klaim dalam praktik, di mana keputusan akhir tetap berada pada auditor, verifikator, dan proses pembuktian yang lebih formal.

\subsection{Keterbatasan Verifikasi Manual pada Volume Klaim Besar}\label{keterbatasan-verifikasi-manual-pada-volume-klaim-besar}

Verifikasi manual memegang peran penting dalam menjaga kesesuaian klaim terhadap aturan koding, aspek klinis, dan kelengkapan administrasi. Panduan verifikasi INA-CBG menunjukkan bahwa proses verifikasi dapat melibatkan perbedaan persepsi antara \emph{coder} fasilitas kesehatan dan verifikator BPJS Kesehatan, terutama ketika dokumen klinis, diagnosis, atau prosedur tidak sepenuhnya sejalan dengan hasil \emph{grouping} klaim \citep{bpjs2018}. Namun, pada saat volume klaim sangat besar dan variabel klaim semakin kompleks, pemeriksaan manual terhadap seluruh klaim menjadi tidak efisien apabila dilakukan tanpa bantuan analitik.

Skala data JKN/BPJS membuat keterbatasan tersebut semakin nyata. Data klaim memuat informasi administratif, diagnosis utama dan sekunder, prosedur, kelas perawatan, histori layanan, hingga komponen tarif yang perlu dibaca secara bersama. Dalam kondisi seperti ini, verifikasi manual membutuhkan mekanisme prioritisasi agar auditor dan verifikator dapat memusatkan perhatian pada klaim yang paling berisiko. Dengan kata lain, persoalan utamanya bukan sekadar menemukan seluruh klaim bermasalah secara otomatis, tetapi menyediakan alat bantu yang mampu menyaring dan mengurutkan klaim berdasarkan indikasi risiko.

\subsection{Kebutuhan Analitik Berbasis Data untuk Prioritas Verifikasi}\label{kebutuhan-analitik-berbasis-data-untuk-prioritas-verifikasi}

Kebutuhan terhadap analitik berbasis data muncul karena data klaim kesehatan bersifat besar, heterogen, dan menyimpan pola yang tidak selalu dapat ditangkap oleh aturan statis. Timofeyev dan Jakovljevic menekankan bahwa \emph{fraud} dan \emph{corruption} dalam \emph{healthcare} menimbulkan kerugian finansial yang besar serta memerlukan respons yang lebih sistematis \citep{timofeyev2022}. Najar dkk. juga menunjukkan bahwa penanganan \emph{medical fraud} dan \emph{abuse} semakin bergeser ke arah integrasi deteksi berbasis data, pencegahan, dan respons hukum \citep{najar2025}. Dalam konteks ini, \emph{machine learning} menjadi relevan bukan karena menggantikan verifikasi, tetapi karena mampu membantu membaca pola dari volume data yang besar dan menandai klaim dengan risiko lebih tinggi.

Atas dasar itu, \emph{problem} yang diangkat dalam penelitian ini dirumuskan sebagai kebutuhan untuk mendeteksi potensi \emph{fraud} atau \emph{upcoding} pada klaim JKN/BPJS secara \emph{claim-centric}, yaitu dengan menjadikan satu klaim sebagai unit analisis utama. Pendekatan ini dibutuhkan agar integrasi atribut peserta, diagnosis, histori layanan, komponen tarif, dan indikator klinis dapat dibaca secara terpadu. Dengan demikian, sistem yang dibangun diarahkan untuk menghasilkan prioritas verifikasi klaim berisiko tinggi secara terukur dan dapat dijelaskan, bukan untuk membuat vonis \emph{fraud} akhir.

\section{TEORI PENUNJANG}\label{teori-penunjang}

Teori penunjang pada penelitian ini disusun untuk menjelaskan dasar konseptual dari domain permasalahan sekaligus dasar metodologis dari solusi yang diusulkan. Pembahasan tidak diarahkan untuk menguraikan implementasi eksperimen secara rinci, melainkan untuk membangun pemahaman mengapa pendekatan \emph{claim-centric}, \emph{pseudo-labeling}, \emph{benchmarking multi-model}, \emph{threshold tuning}, dan interpretabilitas relevan digunakan pada deteksi potensi \emph{fraud} atau \emph{upcoding} dalam sistem klaim JKN/BPJS.

\subsection{\texorpdfstring{\emph{Fraud} Kesehatan dan \emph{Upcoding}}{Fraud Kesehatan dan Upcoding}}\label{fraud-kesehatan-dan-upcoding}

\emph{Fraud} kesehatan secara umum dipahami sebagai tindakan penipuan atau misrepresentasi yang disengaja untuk memperoleh manfaat atau pembayaran yang tidak sah dari sistem pembiayaan kesehatan. Villegas-Ortega dkk. menunjukkan bahwa \emph{fraud} asuransi kesehatan memiliki manifestasi yang beragam, melibatkan \emph{provider}, \emph{beneficiary}, maupun pihak asuransi, serta dipengaruhi oleh faktor regulasi, organisasi, dan perilaku \citep{villegas2021}. Dalam editorial tentang \emph{fraud and corruption in healthcare}, Timofeyev dan Jakovljevic juga menekankan bahwa \emph{fraud} harus dibedakan dari kesalahan yang tidak disengaja karena \emph{fraud} selalu mengandung unsur pelanggaran yang disengaja \citep{timofeyev2022}.

Dalam penelitian ini, bentuk \emph{fraud} atau \emph{abuse} yang menjadi fokus adalah \emph{upcoding}. Crespin dkk. memperlihatkan bahwa pada sistem pembayaran berbasis diagnosis, pengodean pasien ke tingkat intensitas tertinggi dapat meningkat seiring perubahan praktik \emph{coding} dan berpengaruh besar terhadap pembayaran rumah sakit \citep{crespin2024}. Maryana dan Aspan menjelaskan bahwa \emph{upcoding} terjadi ketika diagnosis atau tindakan dikodekan lebih tinggi daripada seharusnya, sehingga menghasilkan nilai klaim yang lebih besar \citep{maryana2024}. Dengan demikian, teori tentang \emph{fraud} kesehatan dan \emph{upcoding} menjadi dasar untuk memahami bahwa ketidakwajaran klaim tidak selalu muncul dari biaya tinggi semata, melainkan dari hubungan antara representasi klinis, aturan \emph{coding}, dan insentif pembayaran.

\subsection{Sistem Klaim JKN/BPJS dan INA-CBG}\label{sistem-klaim-jknbpjs-dan-ina-cbg}

Sistem INA-CBG merupakan mekanisme pembayaran prospektif yang membayar layanan berdasarkan kelompok kasus, bukan berdasarkan setiap tindakan secara terpisah. Permenkes Nomor 27 Tahun 2014 menjelaskan bahwa sistem ini menggunakan prinsip \emph{casemix} dengan pengelompokan diagnosis dan prosedur yang memiliki kemiripan klinis dan pola penggunaan sumber daya \citep{kemenkes2014}. Petunjuk teknis e-claim menjelaskan lebih lanjut bahwa pembentukan klaim berlangsung melalui entri data administratif dan klinis, proses \emph{grouping}, finalisasi, dan pengiriman klaim \citep{kemenkes2016}. Karena itu, kualitas data pada tahap entri dan \emph{coding} berpengaruh langsung terhadap hasil \emph{grouping} dan nilai pembayaran.

Pada level operasional, panduan verifikasi INA-CBG menunjukkan bahwa verifikasi klaim mencakup pemeriksaan administratif, klinis, dan koding \citep{bpjs2018}. Chalkley dkk. memberikan dasar empiris bahwa dalam konteks Indonesia, perubahan tarif dapat berkorelasi dengan perubahan perilaku \emph{coding} rumah sakit \citep{chalkley2022}. Temuan tersebut memperkuat pandangan bahwa sistem pembayaran berbasis kelompok kasus memiliki potensi menciptakan insentif diskresi dalam proses \emph{coding}. Oleh sebab itu, teori tentang JKN/BPJS dan INA-CBG pada penelitian ini tidak hanya menjelaskan sistem pembayarannya, tetapi juga menjelaskan mengapa struktur sistem tersebut relevan terhadap problem deteksi potensi \emph{upcoding}.

\subsection{\texorpdfstring{Data Klaim Kesehatan dan Pendekatan \emph{Claim-Centric}}{Data Klaim Kesehatan dan Pendekatan Claim-Centric}}\label{data-klaim-kesehatan-dan-pendekatan-claim-centric}

Data klaim kesehatan merupakan bentuk data administratif yang merekam interaksi antara peserta, \emph{provider}, diagnosis, prosedur, pelayanan, dan pembiayaan. Dalam penelitian \emph{fraud detection}, data semacam ini bernilai tinggi karena memuat sinyal yang berkaitan dengan pola biaya, utilisasi layanan, hubungan antar kode, serta histori perilaku. Kumaraswamy dkk. menegaskan bahwa \emph{healthcare claims} data dapat ditransformasikan menjadi himpunan fitur sekunder yang lebih informatif untuk mendukung deteksi \emph{fraud} \citep{kumaraswamy2022}. Hamid dkk. dan Nabrawi serta Alanazi juga menunjukkan bahwa data klaim dapat dimanfaatkan untuk membangun pendekatan \emph{data mining} dan \emph{machine learning} pada \emph{healthcare insurance fraud detection} \citep{nabrawi2023}, \citep{hamid2024}.

Pendekatan \emph{claim-centric} berarti satu klaim dijadikan unit observasi utama, lalu diperkaya dengan informasi lain yang relevan. Secara teoretis, pendekatan ini memungkinkan penelitian membaca risiko pada level transaksi klaim secara langsung, bukan pada agregat \emph{provider} atau agregat peserta. Bagi penelitian ini, konsep tersebut penting karena potensi \emph{upcoding} muncul pada klaim tertentu yang dibentuk oleh kombinasi diagnosis, prosedur, \emph{severity level}, tarif, dan histori layanan. Dengan menjadikan klaim sebagai pusat analisis, representasi risiko menjadi lebih dekat dengan objek verifikasi yang sesungguhnya.

\subsection{\texorpdfstring{\emph{Machine Learning} untuk Deteksi Potensi \emph{Fraud}/\emph{Upcoding}}{Machine Learning untuk Deteksi Potensi Fraud/Upcoding}}\label{machine-learning-untuk-deteksi-potensi-fraudupcoding}

\emph{Machine learning} adalah pendekatan komputasional yang memungkinkan sistem belajar dari data untuk mengenali pola, melakukan klasifikasi, atau memprediksi kejadian tertentu. Dalam konteks \emph{fraud detection}, \emph{machine learning} digunakan karena aturan statis sering kali tidak cukup untuk menangkap pola penyimpangan yang kompleks dan berubah. Curtis dkk. menunjukkan bahwa penelitian \emph{healthcare fraud detection} telah memanfaatkan berbagai keluarga \emph{classifier}, baik \emph{supervised}, \emph{unsupervised}, maupun \emph{semi-supervised}, untuk menangani karakter data klaim yang besar, heterogen, dan tidak seimbang \citep{curtis2025}.

Hamid dkk. menekankan bahwa \emph{healthcare insurance fraud detection} menghadapi tantangan dari skema \emph{fraud} yang berkembang, kualitas data yang bervariasi, kebutuhan deteksi yang andal, serta potensi \emph{false positive} yang tinggi \citep{hamid2024}. Nabrawi dan Alanazi juga menunjukkan bahwa model \emph{machine learning} pada \emph{claims} dapat membantu mengenali pola yang sulit ditemukan secara manual \citep{nabrawi2023}. Dalam penelitian ini, teori \emph{machine learning} diposisikan sebagai dasar umum yang menjelaskan mengapa pendekatan klasifikasi berbasis data dapat digunakan untuk membentuk prioritas verifikasi klaim berisiko.

\subsection{\texorpdfstring{\emph{Anomaly Detection}, \emph{Pseudo-Labeling}, dan \emph{Weak/Semi-Supervised Learning}}{Anomaly Detection, Pseudo-Labeling, dan Weak/Semi-Supervised Learning}}\label{anomaly-detection-pseudo-labeling-dan-weaksemi-supervised-learning}

Salah satu keterbatasan mendasar pada \emph{fraud detection} berbasis klaim adalah tidak tersedianya label \emph{fraud} aktual pada level transaksi. Dalam kondisi seperti ini, \emph{anomaly detection} menjadi relevan karena pendekatan tersebut tidak bergantung pada label final. Settipalli dan Gangadharan menunjukkan bahwa pendekatan \emph{unsupervised multivariate analysis} dapat digunakan untuk mendeteksi \emph{fraud} pada \emph{health insurance claims} \citep{settipalli2023}. Shekhar dkk. juga memperlihatkan bahwa \emph{unsupervised machine learning} dapat digunakan untuk \emph{health care fraud detection} sekaligus tetap memberikan penjelasan yang dapat dibaca \citep{shekhar2023}. Meulemeester dkk. lebih lanjut menekankan bahwa \emph{unsupervised anomaly detection} dapat dikembangkan secara \emph{explainable} agar hasilnya bermanfaat untuk investigasi dan audit \citep{meulemeester2025}.

Ketika tujuan penelitian tidak berhenti pada identifikasi anomali, tetapi berlanjut ke model klasifikasi, maka \emph{pseudo-labeling} menjadi strategi yang relevan. \emph{Pseudo-labeling} membentuk target awal dari sinyal anomali atau label lemah yang tersedia, lalu menggunakan target tersebut pada tahap \emph{supervised learning}. Yoon dkk. melalui kerangka SPADE menunjukkan bahwa \emph{semi-supervised anomaly detection} dapat menggunakan \emph{pseudo-labeler} untuk mengatasi keterbatasan label dan \emph{mismatch} distribusi antara data berlabel dan tidak berlabel \citep{yoon2022}. Bakker dkk. juga menunjukkan bahwa kombinasi \emph{pseudo-labeling} dan \emph{anomaly detection} dapat membantu menangani label \emph{scarcity} dalam \emph{billing healthcare} \citep{bakker2025}. Secara teoretis, pendekatan ini sesuai dengan kebutuhan penelitian ini, karena klaim berisiko harus dibentuk dari sinyal awal yang informatif sebelum dibandingkan lebih lanjut melalui model \emph{supervised}.

\subsection{\texorpdfstring{\emph{Feature Engineering} pada \emph{Healthcare Claims}}{Feature Engineering pada Healthcare Claims}}\label{feature-engineering-pada-healthcare-claims}

\emph{Feature engineering} merupakan proses membangun variabel turunan yang lebih representatif daripada atribut mentah agar model dapat menangkap pola yang relevan dengan masalah penelitian. Kumaraswamy dkk. menegaskan bahwa \emph{fraud detection} pada \emph{healthcare claims} sangat bergantung pada kemampuan mengubah data klaim yang heterogen menjadi fitur-fitur sekunder yang mencerminkan \emph{business heuristics}, relasi antar aktor, dan karakter populasi layanan \citep{kumaraswamy2022}. Dengan kata lain, kualitas \emph{feature engineering} sangat menentukan apakah sinyal \emph{fraud} atau \emph{upcoding} dapat muncul secara jelas pada data.

Dalam penelitian ini, teori \emph{feature engineering} penting karena potensi \emph{upcoding} tidak cukup dibaca dari satu variabel mentah seperti tarif atau lama rawat saja. Yang lebih relevan adalah hubungan antar komponen biaya, intensitas diagnosis sekunder, histori layanan, indikator \emph{special} CMG, dan rasio turunan yang memperlihatkan ketidakwajaran pola klaim. Wang dkk. menunjukkan bahwa pemilihan fitur, reduksi fitur, dan interpretasi terhadap fitur yang dominan dapat meningkatkan kualitas prediksi pada \emph{healthcare insurance fraud detection} \citep{wang2025}. Oleh karena itu, \emph{feature engineering} dipahami sebagai proses membangun sinyal risiko klaim yang lebih kaya, bukan sekadar menambah jumlah variabel.

\subsection{\texorpdfstring{\emph{Threshold Tuning} dan Evaluasi pada Data Tidak Seimbang}{Threshold Tuning dan Evaluasi pada Data Tidak Seimbang}}\label{threshold-tuning-dan-evaluasi-pada-data-tidak-seimbang}

\emph{Fraud detection} pada data klaim hampir selalu berhadapan dengan distribusi kelas yang tidak seimbang, karena jumlah klaim berisiko umumnya jauh lebih sedikit dibandingkan klaim normal. Pada kondisi seperti ini, evaluasi tidak cukup mengandalkan satu ukuran umum seperti akurasi. Curtis dkk. dan Hamid dkk. menunjukkan bahwa metrik seperti \emph{precision}, \emph{recall}, \emph{F1-score}, \emph{ROC-AUC}, dan \emph{PR-AUC} lebih relevan karena memberikan gambaran yang lebih tepat mengenai kemampuan model dalam menangkap kelas minoritas dan menjaga kualitas pemeringkatan risiko {[}15, 20{]}. Dalam penelitian ini, \emph{Matthews Correlation Coefficient} juga relevan karena mampu membaca kualitas klasifikasi secara lebih seimbang pada distribusi kelas yang tidak merata.

Selain metrik evaluasi, \emph{threshold tuning} merupakan komponen teoritis yang penting ketika model digunakan untuk prioritas verifikasi. \emph{Threshold} bawaan seperti 0,5 tidak selalu sesuai dengan kebutuhan operasional. Shen dan Kurshan menunjukkan bahwa pemilihan \emph{threshold} pada \emph{fraud alert system} perlu memperhitungkan keseimbangan antara kemampuan deteksi dan kapasitas pemrosesan \emph{alert} \citep{shen2020}. Secara konseptual, hal ini sejalan dengan penelitian ini, karena tujuan model adalah menandai klaim berisiko tinggi untuk diperiksa lebih lanjut, bukan membuat keputusan otomatis final. Oleh karena itu, \emph{threshold tuning} dipahami sebagai proses penyesuaian ambang keputusan model agar lebih sesuai dengan tujuan deteksi risiko pada domain klaim kesehatan.

\subsection{Interpretabilitas Model}\label{interpretabilitas-model}

Interpretabilitas model diperlukan agar keluaran \emph{machine learning} dapat dipahami oleh pengguna domain, khususnya auditor dan verifikator klaim. Pada konteks \emph{fraud detection}, model yang hanya menghasilkan skor tanpa alasan akan sulit dipakai sebagai dasar prioritas pemeriksaan. Salih dkk. menjelaskan bahwa SHAP dan LIME merupakan dua pendekatan yang banyak digunakan dalam \emph{explainable artificial intelligence} untuk memahami perilaku model yang kompleks \citep{salih2025}. SHAP cenderung kuat untuk membaca kontribusi fitur secara konsisten pada tingkat global maupun lokal, sedangkan LIME berguna untuk menjelaskan prediksi individual melalui aproksimasi lokal.

Pada \emph{domain healthcare fraud}, kebutuhan interpretabilitas juga telah ditegaskan oleh beberapa penelitian. Shekhar dkk. dan Meulemeester dkk. menunjukkan bahwa \emph{explainable fraud detection} dapat membantu menghubungkan skor risiko dengan pola yang dapat ditinjau kembali oleh investigator {[}16, 17{]}. Wang dkk. juga menggunakan kombinasi \emph{feature importance}, SHAP, \emph{partial dependence plot}, dan LIME untuk menjelaskan alasan prediksi \emph{fraud} pada data asuransi kesehatan \citep{wang2025}. Dengan demikian, teori interpretabilitas pada penelitian ini menjadi dasar untuk memahami mengapa hasil model harus dapat dibaca baik pada tingkat global, misalnya melalui \emph{feature importance}, SHAP, dan PDP, maupun pada tingkat klaim individual melalui \emph{local explanation} seperti LIME.

\section{PENELITIAN TERKAIT}\label{penelitian-terkait}

Penelitian terkait pada bagian ini dibahas secara komparatif untuk menunjukkan bagaimana studi sebelumnya memandang masalah \emph{fraud} kesehatan, bagaimana pendekatan analitik dibangun, dan di mana letak perbedaan penelitian ini. Ulasan disusun dalam empat klaster utama agar posisi penelitian dapat terlihat lebih jelas, yaitu studi tentang \emph{fraud} atau \emph{upcoding} pada sistem pembayaran berbasis diagnosis, studi fraud JKN di Indonesia, studi \emph{healthcare claims fraud detection} berbasis \emph{machine learning}, serta studi \emph{unsupervised}, \emph{pseudo-labeling}, dan \emph{explainable fraud detection}.

\subsection{\texorpdfstring{Penelitian \emph{Fraud/Upcoding} pada Sistem Pembayaran Berbasis Diagnosis}{Penelitian Fraud/Upcoding pada Sistem Pembayaran Berbasis Diagnosis}}\label{penelitian-fraudupcoding-pada-sistem-pembayaran-berbasis-diagnosis}

Studi pada kelompok ini menekankan bahwa sistem pembayaran berbasis diagnosis dapat menciptakan insentif untuk melakukan diskresi \emph{coding}. Chalkley dkk. menunjukkan bahwa pada sistem DRG di Indonesia terdapat hubungan positif antara perubahan tarif dan perubahan konsentrasi aktivitas berkode tertentu, yang mengindikasikan adanya perilaku \emph{coding} yang sensitif terhadap struktur harga \citep{chalkley2022}. Crespin dkk. memperlihatkan bahwa kenaikan \emph{coding intensity} pada sistem pembayaran berbasis diagnosis dapat berkontribusi besar terhadap pertumbuhan pembayaran rumah sakit. Maryana dan Aspan melengkapi diskusi ini dengan penjelasan yuridis bahwa \emph{upcoding} merupakan perubahan kode diagnosis atau prosedur ke tarif yang lebih tinggi daripada seharusnya \citep{maryana2024}.

Kontribusi kelompok studi ini terletak pada penegasan bahwa \emph{upcoding} bukan sekadar persoalan teknis \emph{coding}, tetapi berkaitan langsung dengan struktur insentif pada sistem pembayaran kesehatan. Namun, studi-studi tersebut umumnya belum membangun \emph{pipeline} analitik berbasis klaim yang terintegrasi sampai tahap pemodelan risiko. Penelitian ini berbeda karena tidak hanya menggunakan temuan tentang insentif \emph{coding} sebagai \emph{problem domain}, tetapi juga mengembangkan mekanisme deteksi berbasis \emph{data claim-centric} untuk mendukung prioritas verifikasi klaim JKN/BPJS.

\subsection{\texorpdfstring{Penelitian \emph{Fraud} JKN/Indonesia}{Penelitian Fraud JKN/Indonesia}}\label{penelitian-fraud-jknindonesia}

Penelitian \emph{fraud} JKN di Indonesia umumnya menyoroti bentuk-bentuk \emph{fraud}, \emph{moral hazard provider}, serta peran rekam medis dan verifikasi dalam menelusuri klaim yang bermasalah. Muliarta dkk. dan Nurlianti dkk. meninjau bahwa \emph{fraud} pada pelaksanaan JKN melibatkan berbagai bentuk penyimpangan dan merupakan persoalan serius dalam tata kelola pembiayaan kesehatan {[}7, 8{]}. Syafrawati dkk. menunjukkan bahwa \emph{moral hazard provider} dalam implementasi asuransi kesehatan sosial dapat muncul dalam bentuk \emph{upcoding} dan perilaku layanan lain yang menyimpang \citep{syafrawati2023}. Sugiarti dkk. menelusuri potensi \emph{fraud} melalui rekam medis dan menunjukkan bahwa ketidaktepatan kode diagnosis serta ketidaksesuaian \emph{clinical pathway} berpengaruh pada nilai klaim \citep{sugiarti2021}.

Kelompok studi ini sangat penting untuk memperkuat konteks lokal Indonesia, tetapi Meioritas masih berfokus pada kajian konseptual, tinjauan literatur, atau audit berbasis dokumen. Pendekatan tersebut memberikan pemahaman domain yang kuat, tetapi belum banyak yang mendorong integrasi data klaim multi-sumber dan pemodelan \emph{machine learning} pada level klaim individual. Di sinilah penelitian ini mengambil posisi yang berbeda, yaitu dengan memanfaatkan \emph{data claim-centric} sebagai objek analisis dan menghubungkannya dengan pendekatan \emph{machine learning} untuk menyusun prioritas verifikasi secara lebih sistematis.

\subsection{\texorpdfstring{Penelitian \emph{Healthcare Claims Fraud Detection} dengan \emph{Machine Learning}}{Penelitian Healthcare Claims Fraud Detection dengan Machine Learning}}\label{penelitian-healthcare-claims-fraud-detection-dengan-machine-learning}

Kelompok berikutnya terdiri atas penelitian yang memanfaatkan \emph{machine learning} untuk mendeteksi \emph{fraud} pada \emph{healthcare claims}. Hamid dkk. menggabungkan \emph{association rule mining} dan teknik \emph{unsupervised learning} untuk mendeteksi \emph{healthcare insurance fraud} \citep{hamid2024}. Nabrawi dan Alanazi menunjukkan bahwa \emph{model machine learning} dapat digunakan untuk meningkatkan deteksi klaim yang mencurigakan pada data asuransi kesehatan \citep{nabrawi2023}. Wang dkk. mengembangkan \emph{model ensemble} yang \emph{robust} dan \emph{interpretable} untuk prediksi \emph{healthcare insurance fraud} dengan dukungan \emph{feature selection} dan teknik interpretabilitas \citep{wang2025}. Curtis dkk. melalui \emph{review} \emph{classifier healthcare fraud detection} menegaskan bahwa berbagai keluarga model telah digunakan, tetapi tantangan penting tetap berkisar pada \emph{class imbalance}, interpretabilitas, kualitas fitur, dan generalisasi model \citep{curtis2025}.

Studi-studi ini menunjukkan bahwa \emph{machine learning} memiliki potensi kuat untuk membaca pola \emph{fraud claims}. Meskipun demikian, sebagian besar penelitian tersebut dibangun pada \emph{dataset} non-Indonesia dan tidak selalu beroperasi pada konteks JKN/BPJS dengan struktur INA-CBG. Selain itu, banyak penelitian telah memiliki \emph{label fraud} yang lebih siap pakai atau menggunakan \emph{setting} yang berbeda dari kondisi klaim JKN. Penelitian ini mengambil posisi berbeda dengan membangun deteksi potensi \emph{fraud} atau \emph{upcoding} pada klaim JKN/BPJS berbasis \emph{claim-centric}, lalu menentukan model final melalui \emph{benchmarking multi-model} secara empiris.

\subsection{\texorpdfstring{Penelitian \emph{Unsupervised}, \emph{Pseudo-Labeling}, dan \emph{Explainable Fraud Detection}}{Penelitian Unsupervised, Pseudo-Labeling, dan Explainable Fraud Detection}}\label{penelitian-unsupervised-pseudo-labeling-dan-explainable-fraud-detection}

Karena label \emph{fraud} aktual pada level klaim tidak selalu tersedia, beberapa penelitian memilih jalur \emph{unsupervised} atau \emph{semi-supervised}. Settipalli dan Gangadharan menggunakan pendekatan \emph{unsupervised multivariate analysis} untuk \emph{fraud detection} pada \emph{health insurance claims} \citep{settipalli2023}. Shekhar dkk. menunjukkan bahwa \emph{unsupervised machine learning} dapat digunakan untuk \emph{explainable healthcare fraud detection}, sedangkan Meulemeester dkk. menekankan pentingnya \emph{explainable unsupervised anomaly detection} untuk data \emph{insurance} {[}16, 17{]}. Yoon dkk. melalui SPADE memberikan dasar metodologis bahwa \emph{semi-supervised anomaly detection} dapat menggunakan \emph{pseudo-labeler} untuk mengatasi keterbatasan label dan \emph{mismatch} distribusi \citep{yoon2022}. Bakker dkk. menunjukkan bahwa \emph{pseudo-labeling} dapat digunakan dalam \emph{billing healthcare} yang mengalami label \emph{scarcity} \citep{bakker2025}. Sementara itu, Kumaraswamy dkk. memperlihatkan bahwa rekayasa fitur yang baik tetap menjadi prasyarat penting agar sinyal risiko dapat muncul secara bermakna \citep{kumaraswamy2022}.

Dari kelompok studi ini terlihat bahwa \emph{anomaly detection}, \emph{pseudo-labeling}, dan \emph{explainability} merupakan kombinasi yang semakin penting dalam deteksi \emph{fraud} berbasis data. Meskipun demikian, masih sedikit penelitian yang secara eksplisit menempatkan klaim JKN/BPJS dalam sistem INA-CBG sebagai unit analisis utama, lalu membangun alur lengkap dari \emph{claim-centric data preparation}, pembentukan \emph{pseudo-label}, \emph{benchmarking multi-model}, \emph{threshold tuning}, dan interpretabilitas global serta lokal. Dengan demikian, celah riset yang diisi oleh penelitian ini terletak pada penggabungan konteks lokal JKN/BPJS dengan \emph{pipeline} analitik \emph{claim-centric} yang utuh, terukur, dan dapat dijelaskan untuk mendukung prioritas verifikasi klaim berisiko tinggi.

Tabel 2.1 menunjukkan bahwa penelitian-penelitian sebelumnya telah memberikan kontribusi penting dalam memahami \emph{fraud} kesehatan, insentif \emph{upcoding}, \emph{fraud detection} berbasis \emph{machine learning}, serta \emph{anomaly detection} yang dapat dijelaskan. Studi-studi tersebut memperlihatkan bahwa problem \emph{fraud} pada layanan kesehatan tidak hanya berkaitan dengan tindakan curang yang bersifat administratif, tetapi juga berhubungan dengan struktur insentif pada sistem pembayaran, kualitas pengkodean diagnosis dan prosedur, serta karakteristik data klaim.

Namun demikian, masih terdapat beberapa celah yang penting. Sebagian besar penelitian masih menggunakan data non-Indonesia, beroperasi pada konteks sistem pembayaran yang berbeda dari JKN/BPJS dan INA-CBG, atau belum membangun alur analitik yang utuh mulai dari pembentukan sinyal risiko awal ketika \emph{label fraud} aktual tidak tersedia sampai penyediaan interpretabilitas hasil pada level klaim. Oleh karena itu, penelitian ini mengambil posisi yang berbeda dengan menempatkan klaim JKN/BPJS dalam sistem INA-CBG sebagai unit analisis utama, lalu membangun pendekatan \emph{claim-centric} yang mencakup \emph{pseudo-labeling}, \emph{benchmarking multi-model}, \emph{threshold tuning}, dan interpretabilitas untuk mendukung prioritas verifikasi klaim berisiko tinggi.

\phantomsection\label{_Toc228109609}{}Tabel 2. 1 Komparasi Penelitian Terkait Deteksi Fraud/Upcoding pada Klaim Kesehatan

\begin{longtable}[]{@{}
  >{\raggedright\arraybackslash}p{(\columnwidth - 14\tabcolsep) * \real{0.0419}}
  >{\raggedright\arraybackslash}p{(\columnwidth - 14\tabcolsep) * \real{0.0845}}
  >{\raggedright\arraybackslash}p{(\columnwidth - 14\tabcolsep) * \real{0.1267}}
  >{\raggedright\arraybackslash}p{(\columnwidth - 14\tabcolsep) * \real{0.1267}}
  >{\raggedright\arraybackslash}p{(\columnwidth - 14\tabcolsep) * \real{0.1795}}
  >{\raggedright\arraybackslash}p{(\columnwidth - 14\tabcolsep) * \real{0.1583}}
  >{\raggedright\arraybackslash}p{(\columnwidth - 14\tabcolsep) * \real{0.1574}}
  >{\raggedright\arraybackslash}p{(\columnwidth - 14\tabcolsep) * \real{0.1251}}@{}}
\toprule\noalign{}
\begin{minipage}[b]{\linewidth}\raggedright
\textbf{No}
\end{minipage} & \begin{minipage}[b]{\linewidth}\raggedright
\textbf{Peneliti (Tahun)}
\end{minipage} & \begin{minipage}[b]{\linewidth}\raggedright
\textbf{Konteks/Data}
\end{minipage} & \begin{minipage}[b]{\linewidth}\raggedright
\textbf{Fokus Utama}
\end{minipage} & \begin{minipage}[b]{\linewidth}\raggedright
\textbf{Pendekatan/Metode}
\end{minipage} & \begin{minipage}[b]{\linewidth}\raggedright
\textbf{Kontribusi Kunci}
\end{minipage} & \begin{minipage}[b]{\linewidth}\raggedright
\textbf{Kelebihan Utama}
\end{minipage} & \begin{minipage}[b]{\linewidth}\raggedright
\textbf{Keterbatasan/Pembeda}
\end{minipage} \\
\midrule\noalign{}
\endhead
\bottomrule\noalign{}
\endlastfoot
1 & Chalkley dkk. (2021) & Data JKN Indonesia, sistem pembayaran berbasis DRG/INA-CBG & Sensitivitas perilaku \emph{coding} rumah sakit terhadap perubahan tarif & Analisis empiris kebijakan pembayaran & Menunjukkan hubungan antara perubahan tarif dan perubahan perilaku \emph{coding} & Sangat relevan dengan konteks Indonesia dan insentif INA-CBG & Tidak membangun model deteksi klaim berisiko pada level klaim individual \\
2 & Crespin dkk. (2024) & \emph{Data discharge-level} rumah sakit pada sistem pembayaran berbasis diagnosis & \emph{Upcoding} pada tingkat intensitas/\emph{severity} tertinggi & Analisis kuantitatif perubahan \emph{coding intensity} & Menunjukkan bahwa peningkatan \emph{coding intensity} dapat berkontribusi besar terhadap pertumbuhan pembayaran & Memperkuat hubungan antara \emph{severity-based payment} dan insentif \emph{upcoding} & Bukan konteks JKN/BPJS dan tidak menawarkan \emph{pipeline machine learning claim-centric} \\
3 & Sugiarti dkk. (2022) & Rekam medis rumah sakit dalam konteks JKN & Penelusuran potensi \emph{fraud} melalui akurasi kode diagnosis dan \emph{clinical pathway} & Studi kuantitatif-kualitatif berbasis dokumen rekam medis & Menunjukkan bahwa ketidaktepatan kode dan \emph{clinical pathway} dapat memengaruhi tarif klaim & Relevan untuk membedakan \emph{upcoding}, \emph{error coding}, dan unsur kesengajaan & Berbasis audit dokumen, belum mengembangkan model analitik berbasis data klaim skala besar \\
4 & Hamid dkk. (2024) & \emph{Healthcare insurance claims} & Deteksi \emph{healthcare insurance fraud} & \emph{Data mining, association rule mining}, dan \emph{unsupervised learning} & Menunjukkan bahwa kombinasi teknik \emph{data mining} dapat mendukung deteksi \emph{fraud} pada \emph{claims data} & Relevan untuk \emph{fraud detection} berbasis data klaim & Tidak berfokus pada konteks JKN/INA-CBG dan tidak menekankan pendekatan \emph{claim-centric Indonesia} \\
5 & Kumaraswamy dkk. (2022) & \emph{Healthcare claims data} & \emph{Feature engineering} untuk \emph{fraud detection} & Rekayasa fitur berbasis \emph{claims} dan \emph{business heuristics} & Menunjukkan pentingnya transformasi \emph{claims data} menjadi fitur sekunder yang lebih informatif & Sangat kuat sebagai dasar teoritis \emph{feature engineering} pada \emph{claims} & Fokus utama pada \emph{feature engineering}, bukan \emph{pipeline} lengkap \emph{pseudo-labeling} sampai interpretabilitas \\
6 & Settipalli \& Gangadharan (2022) & \emph{Health insurance claims} & \emph{Fraud detection} tanpa label fraud aktual & \emph{Unsupervised multivariate analysis} & Menunjukkan bahwa \emph{claims fraud} dapat dideteksi dengan pendekatan \emph{unsupervised} & Relevan untuk dasar \emph{anomaly detection} pada \emph{claims} & Belum menggabungkan \emph{pseudo-labeling}, \emph{benchmarking multi-model}, dan interpretabilitas secara utuh \\
7 & Shekhar dkk. (2022) & \emph{Health care fraud data} & \emph{Explainable fraud detection} tanpa label siap pakai & \emph{Unsupervised machine learning} dan \emph{explainable fraud detection} & Menunjukkan bahwa \emph{fraud detection} dapat dilakukan secara \emph{unsupervised} tetapi tetap \emph{explainable} & Menghubungkan \emph{anomaly detection} dengan interpretabilitas & Bukan konteks JKN/BPJS dan tidak menempatkan klaim INA-CBG sebagai unit analisis utama \\
8 & Meulemeester dkk. (2025) & \emph{Healthcare insurance data} & \emph{Explainable unsupervised anomaly detection} & \emph{Unsupervised anomaly detection} dengan \emph{explainability} & Menunjukkan bahwa \emph{anomaly detection} dapat diarahkan untuk investigasi/audit & Kuat untuk mendukung \emph{explainable anomaly detection} pada insurance data & Tidak membahas konteks INA-CBG dan tidak membangun \emph{supervised benchmarking} setelah \emph{pseudo-labeling} \\
9 & Polak dkk. (2025) & Literatur \emph{healthcare fraud detection} & \emph{Review classifier machine learning} pada \emph{healthcare fraud detection} & \emph{Review} berbagai \emph{classifier supervised}, \emph{unsupervised}, dan \emph{semi-supervised} & Memetakan tantangan utama seperti \emph{class imbalance}, kualitas fitur, interpretabilitas, dan generalisasi & Cocok sebagai payung penelitian terkait ML pada \emph{healthcare fraud} & Bersifat \emph{review}, bukan implementasi pada data JKN/BPJS \\
10 & Yoon dkk. (2022) & \emph{Tabular and image anomaly detection} & \emph{Semi-supervised anomaly detection under distribution mismatch} & SPADE, \emph{pseudo-labeler}, \emph{ensemble one-class classifiers} & Menunjukkan bahwa \emph{pseudo-labeling} dapat membantu ketika label terbatas dan distribusi data tidak seragam & Kuat sebagai dasar metodologis \emph{pseudo-labeling}/\emph{semi-supervised learning} & Bukan studi \emph{healthcare claims} secara langsung \\
11 & Bakker dkk. (2025) & \emph{Mental healthcare provider billing} & \emph{Anomaly detection} pada \emph{billing} dengan \emph{label scarcity} & \emph{Hybrid deep learning, pseudo-labeling via Isolation Forest} dan \emph{Autoencoder} & Menunjukkan bahwa \emph{pseudo-labeling} dapat membantu pada \emph{billing healthcare} dengan label terbatas & Relevan untuk menjelaskan label \emph{scarcity} pada \emph{data billing} kesehatan & Domainnya mental \emph{healthcare billing} dan masih belum spesifik JKN/INA-CBG \\
12 & Penelitian ini (2026) & Klaim JKN/BPJS dalam sistem INA-CBG, unit analisis satu klaim FKRTL & Deteksi potensi \emph{fraud/upcoding} untuk prioritas verifikasi & \emph{Claim-centric data preparation}, \emph{pseudo-labeling}, \emph{branch-based feature selection}, \emph{supervised benchmarking}, \emph{threshold tuning}, dan interpretabilitas & Mengintegrasikan konteks lokal JKN/BPJS dengan \emph{pipeline} analitik \emph{claim-centric} yang utuh & Menghubungkan \emph{problem domain} INA-CBG dengan analitik risiko yang terukur dan dapat dijelaskan & Hasil model tetap diposisikan sebagai indikasi risiko dan prioritas verifikasi, bukan pembuktian \emph{fraud} aktual \\
\end{longtable}

\section*{}

Berdasarkan Tabel 2.1, penelitian-penelitian sebelumnya telah memberikan kontribusi penting dalam memahami fraud kesehatan, insentif upcoding, fraud detection berbasis machine learning, serta anomaly detection yang dapat dijelaskan. Studi-studi tersebut menunjukkan bahwa problem fraud pada layanan kesehatan tidak hanya berkaitan dengan tindakan curang yang bersifat administratif, tetapi juga berhubungan dengan struktur insentif pada sistem pembayaran, kualitas pengkodean diagnosis dan prosedur, serta karakteristik data klaim yang kompleks. Dengan demikian, kajian terdahulu memberikan landasan yang kuat bahwa deteksi potensi fraud/upcoding memang membutuhkan pendekatan berbasis data yang lebih sistematis.

Namun demikian, masih terdapat beberapa celah yang penting. Sebagian besar penelitian masih menggunakan data non-Indonesia, beroperasi pada konteks sistem pembayaran yang berbeda dari JKN/BPJS dan INA-CBG, atau belum membangun alur analitik yang utuh mulai dari pembentukan sinyal risiko awal ketika label fraud aktual tidak tersedia sampai penyediaan interpretabilitas hasil pada level klaim. Di sisi lain, studi-studi lokal Indonesia umumnya masih berfokus pada kajian konseptual, audit dokumen, atau pembahasan fraud dari sisi kebijakan dan verifikasi, sehingga belum banyak yang mengintegrasikan data klaim multi-sumber dan pemodelan machine learning pada level klaim individual.

Oleh karena itu, penelitian ini mengambil posisi yang berbeda dengan menempatkan klaim JKN/BPJS dalam sistem INA-CBG sebagai unit analisis utama, lalu membangun pendekatan claim-centric yang mencakup data preparation, pseudo-labeling, branch-based feature selection, benchmarking multi-model, threshold tuning, dan interpretabilitas global serta lokal. Posisi ini membuat penelitian tidak hanya menerapkan machine learning untuk klasifikasi, tetapi juga merancang mekanisme analitis yang lebih sesuai untuk mendukung prioritisasi verifikasi klaim berisiko tinggi. Dengan kata lain, kontribusi penelitian ini terletak pada penggabungan konteks lokal JKN/BPJS, struktur data berbasis klaim, dan kebutuhan interpretabilitas ke dalam satu alur analisis yang terukur dan dapat dijelaskan.
